% !TeX root = ../navier-stokes-manuscript.tex
\subsection{The whole-terminal estimate and the intersection}\label{sub:BodyCompatibleIntersection}

For each fixed \(T<T_*\), equation \eqref{eq:StrictSubintervalRectangleBound} controls the selected rectangle only up to that cutoff.
The missing quantifier is uniformity in \(T\): after \(N\) and \(M\) are fixed, one constant has to survive every \(T<T_*\).

\begin{theorem}[Datum-generated whole-terminal rectangles]\label{thm:BodyWholeTerminalRectangleTheorem}
Let \(u_0\in\Ssigma\), let \(\nu>0\), and let \((u,p)\) be the maximal classical solution generated by \((u_0,\nu)\) on \([0,T_*)\).
If \(T_*<\infty\), then for every finite \(N,M\in\Nzero\),
\begin{equation}\label{eq:BodyWholeTerminalRectangleEstimate}
  \sup_{0<T<T_*}\mathfrak R_{N,M}(u;T)
  \leq C_{N,M}(u_0,\nu,T_*)<\infty.
\end{equation}
The constant may depend on the selected depths \(N,M\), but it is independent of the terminal cutoff \(T\).
The pressure rows are the rows generated by \eqref{eq:BodyTemporalPressureRecurrence}, and restriction to lower temporal or spatial depths preserves every shared derivative.
\end{theorem}

Fix \(N,M\).
Since the integrands in \eqref{eq:BodyRectangleNorm} are nonnegative, monotone convergence as \(T\uparrow T_*\) gives, for \(0\leq n\leq N\),
\begin{equation}\label{eq:BodyWholeTerminalAdjacentRows}
  V_n\in L^2(I_*;H^{M+1}),
  \qquad
  V_{n+1}=\partial_tV_n\in L^2(I_*;H^{M-1}).
\end{equation}
The corresponding pressure rows satisfy
\begin{equation}\label{eq:BodyWholeTerminalPressureRows}
  Q_n\in L^1(I_*;H^M),
  \qquad 0\leq n\leq N,
\end{equation}
by the pressure recurrence, Riesz-transform boundedness, and finite-order Sobolev product estimates.

Every larger rectangle is cut from the same mixed jet and restricts to every smaller one.
For \(N'\leq N\) and \(M'\leq M\), let \(\rho_{N',M'}^{N,M}\) discard the higher temporal rows and lower the retained Sobolev heights through the natural embeddings.
Then
\begin{equation}\label{eq:BodyRectangleRestrictionCompatibility}
  \rho_{N',M'}^{N,M}\cR_{N,M}[u_0]
  =\cR_{N',M'}[u_0].
\end{equation}
The compatible intersection is the full family of those finite views:
\begin{equation}\label{eq:DatumSpecificCompatibleIntersection}
  \boxed{
  \cI[u_0]
  :=\bigl(\cR_{N,M}[u_0]\bigr)_{N,M\in\Nzero},
  \qquad
  \rho_{N',M'}^{N,M}\cR_{N,M}[u_0]
  =\cR_{N',M'}[u_0].}
\end{equation}
For every finite request, one sufficiently large rectangle supplies the requested coordinate; the compatibility relation identifies it with the same derivative in every larger rectangle.
The restriction law \eqref{eq:BodyRectangleRestrictionCompatibility} and the pressure rows \eqref{eq:BodyWholeTerminalPressureRows} give the complete bonding data used below.

\begin{proposition}[The mixed intersection]\label{prop:BodyCompatibleProjectiveAssembly}
The velocity component of \(\cI[u_0]\) satisfies
\begin{equation}\label{eq:BodyProducedMixedIntersection}
  u\in
  \bigcap_{r=0}^{\infty}
  \bigcap_{m=0}^{\infty}
  H^r(I_*;H^m(\R^3)).
\end{equation}
In particular,
\begin{equation}\label{eq:CompleteSpatialRow}
  u\in\bigcap_{m=0}^{\infty}L^2(I_*;H^m(\R^3)).
\end{equation}
\end{proposition}

\begin{proof}
Fix finite \(r,m\), and choose \(N\geq r\) and \(M\geq m\).
The upper edges of that rectangle give
\[
  \partial_t^nu\in L^2(I_*;H^{M+1})
  \hookrightarrow L^2(I_*;H^m)
  \qquad(0\leq n\leq r).
\]
These are the distributional time derivatives of one velocity field, so \(u\in H^r(I_*;H^m)\).
The pair \((r,m)\) was arbitrary.
Taking \(r=0\) gives the complete spatial row.
\end{proof}

The intersection has placed every fixed temporal row in every finite spatial column on one whole terminal interval.
Section~\ref{sec:TheMissingEstimateAndGlobalSmoothness} now reads concrete estimates from those coordinates and returns them to the critical height and the balanced vortex.
